\documentclass[12pt,a4paper]{article}
\usepackage[margin=1in]{geometry}
\usepackage{graphicx}
\usepackage{booktabs}
\usepackage{array}
\usepackage{xcolor}
\usepackage{hyperref}
\usepackage{listings}
\usepackage{enumitem}
\usepackage{setspace}
\usepackage{float}
\usepackage[T1]{fontenc}
\usepackage[utf8]{inputenc}
\definecolor{codebg}{RGB}{247,249,252}
\definecolor{codeframe}{RGB}{210,219,231}
\definecolor{codeblue}{RGB}{18,59,114}
\lstset{basicstyle=\ttfamily\small,backgroundcolor=\color{codebg},frame=single,rulecolor=\color{codeframe},breaklines=true,showstringspaces=false,numbers=left,numberstyle=\tiny\color{gray},keywordstyle=\color{codeblue}\bfseries,commentstyle=\color{gray},tabsize=2}
\hypersetup{colorlinks=true,linkcolor=codeblue,urlcolor=codeblue}
\setstretch{1.15}
\begin{document}
\begin{titlepage}
\centering
{\Large Southeast University\par}\vspace{0.4cm}
{\large Department of Computer Science and Engineering\par}\vspace{1.5cm}
{\Huge \textbf{Lab Report 2}\par}\vspace{0.5cm}
{\LARGE \textbf{HTML Fundamentals and Structured Webpage Development}\par}\vspace{0.5cm}
{\large SEU Tech Event Information and Registration Page\par}\vfill
\begin{tabular}{>{\bfseries}p{4.2cm} p{8cm}}
Course Code: & CSE 472 \\
Course Title: & Web and Internet Programming Lab \\
Instructor: & Abid Ahmad \\
Student Name: & Rakibul Hasan \\
Student ID: & \textit{[Enter Student ID]} \\
Section: & \textit{[Enter Section]} \\
Semester: & \textit{[Enter Semester]} \\
Submission Date: & 17 September 2026 \\
\end{tabular}\vfill
\end{titlepage}
\tableofcontents\newpage
\section{Objective}
The objective of this lab was to learn and apply the fundamental structure of an HTML5 webpage. The work focused on headings, paragraphs, hyperlinks, images, ordered and unordered lists, nested lists, tables, forms, and semantic HTML elements. Another objective was to organize a complete single-page website using \texttt{header}, \texttt{nav}, \texttt{main}, \texttt{section}, and \texttt{footer} elements. The final task was to create an SEU Tech Event Information and Registration Page and present the work in a professional lab report.
\section{Software and Tools Used}
\begin{itemize}[leftmargin=1.5em]
\item Visual Studio Code or any plain-text code editor
\item Google Chrome / Microsoft Edge / Mozilla Firefox
\item HTML5 for webpage structure
\item CSS3 for page presentation and responsive layout
\item Git and GitHub for version control and submission
\end{itemize}
\section{Short Theory}
HTML stands for HyperText Markup Language. It is a markup language used to define the structure and meaning of webpage content. A browser reads the HTML source code and renders it as a visible webpage. A correct HTML5 document normally begins with \texttt{<!DOCTYPE html>} and contains the root \texttt{html} element with separate \texttt{head} and \texttt{body} sections.

The \texttt{head} section stores metadata such as character encoding, viewport settings, page title, and external resource links. The visible page content is written inside \texttt{body}. Headings and paragraphs create content hierarchy, while anchor and image tags add links and visual content. Lists organize related items, tables display data in rows and columns, and forms collect user input. Semantic elements such as \texttt{header}, \texttt{nav}, \texttt{main}, \texttt{section}, and \texttt{footer} make the webpage easier to understand and maintain.
\section{Task Description}
The assigned webpage is titled \textbf{``SEU Tech Event Information and Registration Page''}. The page includes a proper HTML5 document structure, semantic layout, a working external hyperlink, an image with alternative text, ordered and unordered lists, a nested list, a schedule table with more than four data rows, and a registration form containing name, student ID, email, preferred date, radio buttons, checkboxes, and submit/reset buttons.
\section{Working Steps}
\begin{enumerate}[leftmargin=1.8em]
\item Created a project folder named \texttt{lab-02-html}.
\item Created the main file \texttt{index.html} and added the standard HTML5 skeleton.
\item Added metadata, title, and a stylesheet link inside the \texttt{head} section.
\item Built a header containing the event title and introductory text.
\item Added a navigation bar with internal links to About, Topics, Schedule, and Register sections.
\item Created the About section with headings, paragraphs, a meaningful \texttt{div}/\texttt{span}, an external hyperlink, and a banner image.
\item Added unordered, ordered, and nested lists to present event goals, participation steps, and technology tracks.
\item Created an event schedule using a table with a header row and five data rows.
\item Built a registration form with labels, text input, email input, date input, radio buttons, checkboxes, and buttons.
\item Added a footer with contact and course information.
\item Created \texttt{style.css} to improve readability, spacing, responsiveness, navigation, table styling, and form presentation.
\item Tested the page in a browser and checked links, layout, input labels, table structure, and image paths.
\item Uploaded the final code to GitHub under \texttt{labs/lab-02-html}.
\end{enumerate}
\section{Explanation of Important Tags Used}
\begin{table}[H]\centering\small
\begin{tabular}{p{3.2cm} p{10.3cm}}
\toprule\textbf{Tag / Element} & \textbf{Purpose in the Webpage} \\
\midrule
\texttt{<!DOCTYPE html>} & Declares an HTML5 document. \\
\texttt{html} & Root element of the complete webpage. \\
\texttt{head} & Stores metadata, title, viewport settings, and stylesheet reference. \\
\texttt{body} & Contains all visible webpage content. \\
\texttt{header} & Holds the introductory event title and event summary. \\
\texttt{nav} & Contains internal navigation links. \\
\texttt{main} & Wraps the primary content of the page. \\
\texttt{section} & Separates meaningful content areas. \\
\texttt{h1--h3} & Creates a clear heading hierarchy. \\
\texttt{p} & Displays descriptive paragraphs. \\
\texttt{div} & Groups related block-level content. \\
\texttt{span} & Marks a small inline part of text for emphasis. \\
\texttt{a} & Creates internal and external hyperlinks. \\
\texttt{img} & Displays the event banner with alternative text. \\
\texttt{ul}, \texttt{ol}, \texttt{li} & Create unordered, ordered, and nested lists. \\
\texttt{table}, \texttt{tr}, \texttt{th}, \texttt{td} & Display the event schedule. \\
\texttt{form} & Contains the registration interface. \\
\texttt{label} & Describes and connects text with a form control. \\
\texttt{input} & Creates text, email, date, radio, and checkbox controls. \\
\texttt{button} & Creates submit and reset controls. \\
\texttt{footer} & Shows closing contact and course information. \\
\bottomrule
\end{tabular}\caption{Important HTML elements used in Lab Report 2.}
\end{table}
\section{Final Source Code}
\subsection{HTML File: index.html}
\lstinputlisting[language=HTML]{index.html}
\subsection{CSS File: style.css}
\lstinputlisting[language=CSS]{style.css}
\section{Output Screenshots}
The final report should contain at least two clear screenshots of the completed webpage. Save them in a \texttt{screenshots} folder as \texttt{output-top.png} and \texttt{output-bottom.png}.
\begin{figure}[H]\centering
\IfFileExists{screenshots/output-top.png}{\includegraphics[width=0.95\textwidth]{screenshots/output-top.png}}{\fbox{\parbox[c][5cm][c]{0.9\textwidth}{\centering Insert \texttt{screenshots/output-top.png} here.\\Capture the header, navigation, About section, and event image.}}}
\caption{Top section of the completed SEU Tech Event webpage.}\end{figure}
\begin{figure}[H]\centering
\IfFileExists{screenshots/output-bottom.png}{\includegraphics[width=0.95\textwidth]{screenshots/output-bottom.png}}{\fbox{\parbox[c][5cm][c]{0.9\textwidth}{\centering Insert \texttt{screenshots/output-bottom.png} here.\\Capture the schedule table and registration form.}}}
\caption{Schedule and registration area of the completed webpage.}\end{figure}
\section{Problems Faced and Solutions}
\begin{table}[H]\centering\small
\begin{tabular}{p{5.2cm} p{8.3cm}}
\toprule\textbf{Problem} & \textbf{Solution} \\
\midrule
Image not appearing & Verified the \texttt{src} path to \texttt{images/tech-event-banner.svg}. \\
Navigation link not moving to a section & Matched each navigation \texttt{href} value with the correct section \texttt{id}. \\
Table becoming too wide on smaller screens & Used a horizontally scrollable table wrapper. \\
Form labels and controls looking unorganized & Used grouped fields, CSS Grid, consistent spacing, and linked labels. \\
Page becoming crowded on mobile devices & Added responsive media queries for a one-column layout. \\
\bottomrule
\end{tabular}\caption{Problems encountered and their solutions.}
\end{table}
\section{Result and Discussion}
The completed page demonstrates the main HTML topics required in the lab. It has a correct document structure, semantic page sections, headings, paragraphs, an external link, an image, three types of list usage, a structured table, and a complete registration form. CSS improves the presentation while keeping the required HTML concepts clear.
\section{Conclusion}
This lab provided practical experience in building a complete webpage using HTML fundamentals and semantic structure. The final SEU Tech Event Information and Registration Page combines text, links, images, lists, tables, and forms in one organized document. The exercise reinforced proper indentation, meaningful tag selection, file organization, accessibility attributes, and browser testing.
\end{document}
